\section{Task-Information Audit}
\label{sec:method}

A graph intervention changes wiring; its interpretation depends on the
computation that remains. We audit an intervention at three levels. First, we
verify the structural properties the edit claims to preserve. Second, we
examine which task-relevant source signals reach the readout under the
executed operator. Third, we evaluate how a finite trained model uses the
resulting representation. These levels answer different questions and are not
interchangeable evidence.

\subsection{Specify the executed computation}
\label{sec:operator}

For a fixed linear propagation prefix, source-to-readout transfer can be
computed directly, without training a classifier or enumerating every source
input. The audit takes four inputs: an edited graph, its executed propagation
operator, designated task sources, and readout locations.
Let $A_{uv}=1$ denote a directed edge $u\!\to\!v$. The main experiments
use binary support and incoming-neighbor averaging,
\begin{equation}
 d_v=\sum_u A_{uv},\qquad
 W_{vu}=\frac{A_{uv}}{\max(d_v,1)}.
 \label{eq:operator}
\end{equation}
Rows for zero in-degree remain zero. Features propagate linearly for two
rounds, $H=W^2X$, before an MLP reads the concatenated features at
$S=\{s_1,s_2\}$. The operator direction, normalization, and readout
locations are part of the computational model.

With nonzero features only at sources $a,b$, let $E_{ab}$ place the source
vectors and $R_S$ select the sinks. The graph affects the classifier input
through
\begin{equation}
 H_S=T(A)X_{ab},\qquad T(A)=R_SW(A)^2E_{ab}.
 \label{eq:transfer}
\end{equation}
$T$ has two source columns and two sink rows. It records every two-step
route contributing to these inputs, including routes outside a selected
motif family. Equal $T$ matrices give identical readout inputs for any
source features; equal graph statistics need not. Computing $T$ requires
propagating two basis columns, rather than constructing all entries of
$W^2$. For a sparse fixed operator and $h$ steps, this costs
$O(h|E|\,n_{\rm source})$ time and $O(Nn_{\rm source})$ storage.

For this operator, every entry has an explicit path interpretation:
\begin{equation}
 T_{sj}=\sum_r W_{sr}W_{rj}
 =\sum_r\frac{A_{jr}A_{rs}}{\max(d_r,1)\max(d_s,1)}.
 \label{eq:pathsum}
\end{equation}
Each term is one source--relay--sink route. A common-relay count records
only routes shared by both sources. Equation~\eqref{eq:pathsum} also retains
separate-relay routes and their weights. This is why equal degree, module
and common-relay statistics need not imply equal readout inputs.

\subsection{Read the transfer at three levels}
\label{sec:observability}

\paragraph{Source access is not source recovery.}
Source visibility is weaker than source recovery, and neither alone determines an arbitrary trained classifier. For the noiseless XOR task considered here, visibility is nevertheless an exact criterion for whether the label remains recoverable. A source is \emph{visible} when its column of $T$ is nonzero. Visibility
means that some source-dependent signal reaches the readout; it does not
mean that the source vector can be reconstructed. Full column rank is
the stronger linear-separability condition for independent sources.

\paragraph{Ideal task information.}
For the archived subspace task, visibility is already an exact criterion.
Let $c_a,c_b$ be independent balanced bits and
$x_a=P_{c_a}z_a$, $x_b=P_{c_b}z_b$, with independent
$z_a,z_b\sim\mathcal N(0,I_8)$. The class subspaces
$U_c=\operatorname{span}(P_c)\subset\mathbb R^{16}$ satisfy
$U_0\cap U_1=\{0\}$, and the label is $y=c_a\oplus c_b$.

\noindent\textbf{Proposition 1 (ideal population criterion).}
\emph{For a fixed nonnegative $T$ and no other active inputs, the Bayes
accuracy from $H_S=TX_{ab}$ is}
\begin{equation}
 \operatorname{Acc}^{*}(T)=
 \begin{cases}
  1, & T_{:a}\ne0\ \text{and}\ T_{:b}\ne0,\\
  \tfrac12, & \text{otherwise}.
 \end{cases}
 \label{eq:criterion}
\end{equation}

\emph{Proof.} If $T_{:b}=0$, the readout is independent of $c_b$,
so $\Pr(y=1\mid H_S)=1/2$; the same holds for $a$.
If both columns are nonzero, summing sink vectors gives
$z=\alpha x_a+\beta x_b$, where
$\alpha=\sum_sT_{sa}>0$ and $\beta=\sum_sT_{sb}>0$.
For equal source classes, $z\in U_0\cup U_1$. For different classes,
transversality and nonzero Gaussian coordinates imply
$z\notin U_0\cup U_1$ almost surely. Subspace membership therefore
recovers the XOR label. $\square$

Rank-one transfer can thus retain ideal task information: two nonzero
parallel columns still supply both class-dependent terms. A common relay
is likewise unnecessary because separate routes can expose both sources.
The supplement gives the full argument and ten-seed basis-rank check.
The proposition is a population statement for this noiseless task, not a
prediction of finite-sample optimization or a guarantee of source recovery.

\paragraph{Multi-source mixing geometry.}
With $m$ designated active sources, the same linear prefix gives
$H_S=TX$: each column of $T$ belongs to one source and each row of $X$
holds that source's features. For a target source $j$, let $t_j=T_{:j}$
and let $T_{-j}$ collect the other columns. Its private component is the
residual after projecting out those columns:
\begin{equation}
 q_j=(I-P_{\operatorname{span}(T_{-j})})t_j,\qquad
 \rho_j=\|q_j\|_2^2.
 \label{eq:private}
\end{equation}
Because $q_j$ is orthogonal to every other column,
$q_j^\top H_S=\rho_j x_j$ under noiseless mixing. If $\rho_j>0$,
$x_j=q_j^\top H_S/\rho_j$ recovers the target without reconstructing the
other sources. A zero residual means no linear projection isolates it from
arbitrary competing features. That is source recovery, not label solvability:
class information can remain when exact separation is impossible.

\subsection{What the structural interventions constrain}
\label{sec:paths}

The graph controls preserve each node's in-/out-degree and module-pair
edge counts. For fixed sources and sinks, a common-relay family is
\begin{equation}
 \mathcal C(A)=\{(r,s):s\in S,\ A_{ar}A_{br}A_{rs}=1\}.
 \label{eq:paths}
\end{equation}
A member represents two source edges converging at a relay and one edge
to a sink. Its three-edge description omits separate-relay routes.
The number of named motif edges retained is an even narrower statistic.
Neither quantity specifies all columns of Eq.~\eqref{eq:transfer}.

We compare two archived intervention designs. Named-edge compensation
removes selected motif edges and replaces each through a module-matched
double-edge swap
$(u\!\to\!v,x\!\to\!y)\mapsto(u\!\to\!y,x\!\to\!v)$.
Matching source-module and target-module labels preserves module counts;
the exchange preserves every labeled degree. It does not constrain
other source-specific routes.

The path-conditioned design retains a prefix of $\mathcal C$, protects its
path-edge union and rejects swaps that create new members. Other edges
are processed in sorted order with randomized feasible partners. This is
a common-relay constraint, not a complete source-transport constraint.

Degree nulls rewire the whole dose graph. Protected nulls retain selected
path edges while rewiring the background without the no-new-relay rule.
The two therefore ask different counterfactual questions.

\subsection{Turn the analysis into an audit}
\label{sec:controls}

The procedure has four steps. First, verify the declared structural
invariants. Second, construct $T$ from the actual direction, normalization,
propagation depth and sinks. Third, report visibility, the applicable
ideal-task criterion, and, for competing sources, rank and private
components. Finally, compare finite readouts against these information-level
predictions. An oracle given both sources tests label learnability;
$W=I$ prevents cross-node transport; linear-readout, shuffled-label and
moved-readout controls test alternative explanations.

The output pairs a structural certificate with a statement about task
information and a separate measurement of finite-model behavior.

The audit does not replace structural controls. It tests whether those
controls support the task-level conclusion drawn from them. The exact matrix
reduction applies to a fixed linear propagation prefix. Nonlinear message
passing, gating and spiking dynamics require an input-dependent analysis of
source access.
